% !TeX root = ../../Book1.tex
% 纲要标题：### 16.1 可分多项式
% 纲要位置：计划纲要.md，第 558 行。

\section{可分多项式}
\label{b1:sec:1601}


% 纲要内容（仅作写作计划，尚非教材正文）：
%
% * 重根、形式导数与最大公因子
% * 可分元与可分扩张
% * 特征 p 的不可分现象；区分不同根数与根的共同重数
%

扩张次数是相应向量空间的维数。上一章的\cref{b1:lem:simple-embedding-extension}说明：给定基域嵌入，若目标域包含系数变换后的最小多项式的全部根，则单扩张的延拓数等于这些根的不同个数。它何时等于扩张次数，取决于最小多项式是否有重根。特征零时不可约性已经排除了重根；正特征中导数可能恒为零，这会产生新的现象。本节始终在代数闭包中比较根。

\subsection{重根、形式导数与最大公因子}
\label{b1:subsec:1601:01}

\begin{definition}\label{b1:def:separable-polynomial}
设 \(K\) 为域，\(f\in K[x]\) 为非零多项式。如果 \(f\) 在 \(K\) 的一个代数闭包中没有重根，则称 \(f\) 为\term[kefen-duoxiangshi]{可分多项式}[separable polynomial]。这一条件与代数闭包的选择无关；非零常数多项式按此约定也可分。
\end{definition}
次数为 \(n\) 的非零多项式在代数闭包中按重数共有 \(n\) 个根，所以可分性等价于它有 \(\deg f\) 个不同根；非零常数多项式的次数和根数都为零。由\cref{b1:prop:algebraic-closure-unique}，任意两个 \(K\) 的代数闭包之间存在保持 \(K\) 的同构。将根的分解沿此同构搬过去，各根及其重数都被保持，这就说明定义与代数闭包的选择无关。

\ProseParagraphBreak
\cref{b1:prop:multiple-root}的重根判别给出 \(f\) 可分当且仅当 \(\gcd(f,f')=1\)。虽然重根是在代数闭包中检查的，最大公因子却只需在 \(K[x]\) 中计算。若它为一，就有 \(u,v\in K[x]\) 使
\[
uf+vf'=1.
\]
这条等式进入任意扩域后仍成立；若 \(f,f'\) 有公共根，代入就会得到 \(0=1\)。反之，若 \(K[x]\) 中的最大公因子非常数，它在代数闭包中必有根；由于它同时整除 \(f,f'\)，这个根便是二者的公共根。

\begin{proposition}\label{b1:prop:irreducible-separable-derivative}
设 \(K\) 为域，\(f\in K[x]\) 为不可约多项式。则 \(f\) 可分当且仅当 \(f'\neq0\)。因此特征零域上的所有不可约多项式均可分。
\end{proposition}
\begin{proof}
记首一最大公因子为 \(d=\gcd(f,f')\)。它整除 \(f\)，而 \(f\) 不可约，所以 \(d\) 只能为一或与 \(f\) 相伴；这正是不可约性在证明中的作用。若 \(f'\neq0\)，则 \(\deg f'<\deg f\)，第二种可能会要求 \(f\mid f'\)，与次数不等式矛盾。因此 \(d=1\)，\(f\) 可分。若 \(f'=0\)，最大公因子就是首一化的 \(f\)，非常数，故不可分。特征零时次数为 \(n\geq1\)、首项系数为 \(a_n\) 的多项式，其导数首项系数 \(na_n\neq0\)，所以不可约多项式总属于前一种情形。
\end{proof}
特征零保证的是不可约多项式全部可分；一般多项式的不可约分解中仍可能重复出现同一个因子，例如 \((x-1)^2\)。不同的不可约因子互素，由 Bézout 关系知它们没有公共根，所以在特征零域上，一个非零多项式可分恰好等价于其不可约分解中没有重复因子。

\subsection{可分元与可分扩张}
\label{b1:subsec:1601:02}

\begin{definition}\label{b1:def:separable-extension}
设 \(L/K\) 为域扩张，\(a\in L\) 在 \(K\) 上代数。如果最小多项式 \(m_{a,K}\) 可分，则称 \(a\) 为 \(K\) 上的\term[kefen-yuan]{可分元}[separable element]，也说它在 \(K\) 上可分。如果 \(L/K\) 是代数扩张，并且 \(L\) 的每个元素都在 \(K\) 上可分，则称 \(L/K\) 为\term[kefen-kuozhang]{可分扩张}[separable extension]。
\end{definition}
若 \(a\) 被一个可分 \(K\) 系数多项式 \(f\) 零化，其最小多项式 \(m_{a,K}\) 整除 \(f\)。若这个因子有重根，该根在 \(f\) 中也至少出现两次，便与 \(f\) 可分矛盾。因此 \(m_{a,K}\) 可分，\(a\) 也可分。

\ProseParagraphBreak
若扩大基域为中间域 \(K\subseteq E\subseteq L\)，原来的关系 \(m_{a,K}(a)=0\) 仍成立，所以 \(m_{a,E}\) 在 \(E[x]\) 中整除 \(m_{a,K}\)。增加可用的系数可能让原最小多项式分解，并使新的最小多项式次数降低；但新的多项式只是原多项式的因子，不能产生新的重根。故 \(a\) 在 \(K\) 上可分就仍在 \(E\) 上可分。反方向一般不成立：在 \(K(a)\) 上，\(a\) 的最小多项式总是一次的，即使它在 \(K\) 上不可分。

\ProseParagraphBreak
例如 \(\mathbb Q(\sqrt2)/\mathbb Q\) 既可分又正规：特征零给出可分性，而它是 \(x^2-2\) 的分裂域。\(\mathbb Q(\sqrt[3]2)/\mathbb Q\) 同样可分，却只包含 \(x^3-2\) 的实根，因而不正规。可分性检查每个最小多项式的根是否相异；正规性检查这些根是否全部留在所讨论的域里。正特征中还会出现正规但不可分以及两者都不满足的扩张，下面的算例将直接检验它们。

\subsection{特征 \texorpdfstring{\(p\)}{p} 的不可分现象}
\label{b1:subsec:1601:03}

设 \(\operatorname{char}K=p>0\)。由导数公式，\(f'=0\) 当且仅当 \(f\) 的所有非零幂指数均被 \(p\) 整除，即 \(f(x)=g(x^p)\)。这并不必然使 \(f\) 在 \(K[x]\) 中可约，因为系数未必有 \(p\) 次根。

\begin{proposition}\label{b1:prop:inseparable-polynomial-shape}
设 \(K\) 为特征 \(p>0\) 的域，\(f\in K[x]\) 为首一不可约多项式。则存在非负整数 \(r\) 及首一不可约可分多项式 \(g\in K[x]\)，使
\[
f(x)=g(x^{p^r}).
\]
\(f\) 的每个不同根的重数均为 \(p^r\)，不同根数为 \(\deg g\)。若 \(a\) 是 \(f\) 的根，则 \(g\) 是 \(a^{p^r}\) 在 \(K\) 上的最小多项式，因而 \(a^{p^r}\) 在 \(K\) 上可分。
\end{proposition}
\begin{proof}
只要导数为零，就把所有非零项的正指数同时除以 \(p\)。多项式仍非常数而次数严格下降，所以有限步后得到导数非零的 \(g\)。这里 \(r\) 是使 \(f\) 的所有非零项指数都被 \(p^r\) 整除的最大整数：所得 \(g\) 至少有一个非零项的指数不被 \(p\) 整除，否则还可以再除一次。若 \(g\) 可约，把分解中的变量换成 \(x^{p^r}\) 就使 \(f\) 可约，矛盾。故 \(g\) 不可约，由\cref{b1:prop:irreducible-separable-derivative}可分。

在代数闭包中，设 \(g\) 的互异根为 \(b_1,\ldots,b_s\)。每个 \(b_i\) 有唯一 \(p^r\) 次根 \(a_i\)：存在性来自代数闭性；若 \(u^{p^r}=v^{p^r}\)，则 \((u-v)^{p^r}=0\)，而域没有非零幂零元，所以 \(u=v\)。这正是幂映射 \(u\mapsto u^p\) 及其 \(r\) 次迭代的单射性。于是
\[
f(x)=\prod_{i=1}^s(x^{p^r}-b_i)=\prod_{i=1}^s(x-a_i)^{p^r}.
\]
不同 \(b_i\) 对应不同 \(a_i\)，故不同根数为 \(s=\deg g\)，每个根的重数为 \(p^r\)，并且
\[
\deg f=(\deg g)p^r.
\]
因此 \(g\) 的次数记录了全部不同根的个数，而 \(p^r\) 记录它们共有的重数。又因为 \(a^{p^r}\) 是首一不可约多项式 \(g\) 的根，\(g\) 就是它在 \(K\) 上的最小多项式，且 \(a^{p^r}\) 在 \(K\) 上可分。
\end{proof}
在 \(K=\mathbb F_p(t)\) 上，\(x^p-t\) 不可约：它在唯一分解整环 \(\mathbb F_p[t]\) 上对素元 \(t\) 满足 Eisenstein 判据，再由 Gauss 分解定理得到分式域上的不可约性。在一个代数闭包中取 \(u\) 使 \(u^p=t\)，并令 \(L=K(u)\)。此时 \([L:K]=p\)，但在 \(L[x]\) 中有
\[
x^p-t=(x-u)^p.
\]
这个等式给出一个重数为 \(p\) 的根，故 \(u\) 在 \(K\) 上不可分，\(L/K\) 也不可分；同时 \(L\) 已是该多项式的分裂域，由\cref{b1:thm:normal-embedding-criterion}知 \(L/K\) 正规。不可约性是在 \(K[x]\) 中成立的，进入 \(L[x]\) 后同一个多项式已经分解，所以不可约与否必须指明系数域。

\ProseParagraphBreak
正规性和可分性也可以同时失败。取 \(K=\mathbb F_2(t)\)，在代数闭包中取 \(a^6=t\)，令 \(L=K(a)=\mathbb F_2(a)\)。多项式 \(x^6-t\) 同样对素元 \(t\) 满足 Eisenstein 判据，故在 \(K[x]\) 中不可约；其导数为零，所以 \(L/K\) 不可分。在代数闭包中有
\[
x^6-t=(x^3-a^3)^2
=(x-a)^2(x^2+ax+a^2)^2.
\]
若它在 \(L\) 中完全分裂，\(L\) 就包含一个满足 \(\zeta^2+\zeta+1=0\) 的三次单位根 \(\zeta\)。可是 \(a\) 超越于 \(\mathbb F_2\)；若写 \(\zeta=P(a)/Q(a)\)，其中非零 \(P,Q\in\mathbb F_2[x]\)，就会得到 \(P^2+PQ+Q^2=0\)。当 \(\deg P\ne\deg Q\) 时，平方次数较高的一项无法消去；当次数相同时，最高项系数为 \(1+1+1=1\)，仍不能消去。因此这样的 \(\zeta\) 不在 \(L\) 中，\(x^6-t\) 不在 \(L\) 中分裂，\(L/K\) 也不正规。
